\chapter{Introdução}
\label{chap:introducao}

\section{Enquadramento e motivação}

A preparação de uma equipa de basquetebol exige decisões sobre o cinco
inicial, a distribuição de minutos e as rotações, conciliando o desempenho
imediato com a disponibilidade do plantel para jogos posteriores. Os dados
históricos, as estatísticas individuais e os registos \emph{play-by-play}
permitem analisar resultados e alinhamentos, apoiando a comparação de
alternativas pelo treinador.

Desenvolvido nas unidades curriculares de Aprendizagem Automática (AAUTIA)
e Planeamento em Inteligência Artificial (PLANTIA), este trabalho propõe
um sistema inteligente de apoio à decisão que articula previsões de
resultados e desempenho com a seleção de ações sujeitas a restrições.

\section{Problema de investigação}

Prever a probabilidade de vitória não basta para recomendar decisões:
é necessário representar os jogadores elegíveis, os alinhamentos, os
minutos de utilização e os objetivos do treinador. As previsões devem
ser compatíveis com as alternativas avaliadas pelo planeador, pois
estatísticas agregadas podem ser insuficientes para comparar alterações
concretas na equipa.

A questão central é: \emph{em que medida a integração de previsões de
aprendizagem automática com métodos de planeamento e otimização pode
melhorar as decisões simuladas de gestão de uma equipa de basquetebol,
face a estratégias fixas ou heurísticas simples?}

A avaliação distinguirá a qualidade preditiva da qualidade das decisões.
As associações históricas não demonstram efeitos causais de substituições
ou alterações de minutos; as conclusões dependerão das hipóteses e da
validade do ambiente de simulação.

\section{Objetivos do trabalho}

O objetivo geral é desenvolver e avaliar um protótipo que integre previsão
e planeamento. Os objetivos específicos são:

\begin{enumerate}
  \item Organizar dados históricos e construir indicadores com informação
  disponível no momento de cada previsão.
  \item Comparar modelos supervisionados para prever vitórias e, quando
  possível, o desempenho de alinhamentos.
  \item Analisar os fatores que contribuem para as previsões e facilitar
  a sua interpretação.
  \item Formular a seleção de jogadores e a distribuição de minutos sob
  restrições de disponibilidade, utilização e composição da equipa.
  \item Comparar métodos de otimização e planeamento num ambiente comum,
  incluindo sequências de vários jogos.
  \item Avaliar a integração dos módulos quanto à previsão, ao cumprimento
  das restrições e aos resultados simulados.
\end{enumerate}

\section{Abordagem proposta e âmbito}

A componente de AAUTIA abrangerá a preparação dos dados, a construção de
características e a comparação de modelos como Regressão Logística,
Random Forest e XGBoost. A avaliação respeitará a ordem temporal dos jogos,
evitando informação futura, e analisará a classificação e a qualidade das
probabilidades utilizadas pelo planeador.

A componente de PLANTIA representará estados, ações, restrições e objetivos
com base nas estimativas disponíveis. Serão consideradas a otimização de
alinhamentos e o planeamento por simulação ou processos de decisão de
Markov. As estratégias serão comparadas com referências simples, sob as
mesmas condições iniciais e restrições.

O âmbito inicial será a previsão pré-jogo e a avaliação em simulação,
condicionadas pela disponibilidade e qualidade dos dados. A fadiga, se
incluída, exigirá hipóteses explícitas de modelação. A previsão durante o
jogo e a aprendizagem por reforço serão possíveis extensões. O protótipo
apoiará a análise do treinador, sem pressupor ganhos em jogos reais ou
resultados experimentais já demonstrados.

\section{Organização do relatório}

Após a introdução, o Capítulo~\ref{chap:estado-arte} apresenta o estado da
arte e a análise crítica que fundamenta a proposta.
O Capítulo~\ref{chap:data} aborda a seleção e preparação dos dados;
o Capítulo~\ref{chap:metodologia}, os modelos e o planeamento; o Capítulo~\ref{chap:experiencias}, o protocolo
experimental, os resultados e a discussão.
O Capítulo~\ref{chap:conclusoes} reunirá as conclusões, as limitações e o
trabalho futuro.
